%% ============================================================
%% appendix_mergechanges.tex — 부록 D: 중간 점검 머지 — 주요 코드 변경
%% 부록 A(1일차 코드 분석, app:code)가 기록한 머지 전 코드 상태와,
%% 4~5일차 브랜치 통합 머지로 반영된 현재 main 을 대조 정리한다.
%% 근거: git diff 6f49e27(통합 전 main) .. 통합 후 main
%% ============================================================
\section{중간 점검 머지 --- 주요 코드 변경 (부록 A 대비)}
\label{app:merge}

본 부록은 \textbf{5일차 작업 브랜치 통합 머지}(2장 5일차 참조)로 \code{main}에 반영된
주요 코드 변경을, \textbf{부록~\ref{app:code}(1일차 코드 분석 종합 보고서)가 기록한 머지 전
코드 상태}와 대조해 정리한 \emph{중간 점검}이다. 부록 A가 \emph{통합 전} 코드를 기술한
기준선이라면, 본 부록은 그 위에 1$\cdot$2차 Refactoring 계획(부록~\ref{app:refactor1}$\cdot$%
\ref{app:refactor2})이 \emph{실제로 실행}되어 달라진 점을 요약한다. 근거는 통합 전
\code{main}과 통합 후 \code{main}의 형상 비교(\code{git diff})이다.

\subsection{변경 규모 요약}
통합 머지로 코드 도메인(\code{Web/}$\cdot$\code{Pico/}$\cdot$\code{AI/})에서 35개 파일이
바뀌었다(\textbf{+5,152 / $-$2,563} 라인, 메시 자산$\cdot$\code{simulation\_backup.js} 제외).
핵심 파일의 전후 비교는 다음과 같다.

\begin{table}[H]\centering\small
\caption{핵심 파일 전후 비교 --- 부록~\ref{app:code} 기준선 대비 통합 후 \code{main}}
\renewcommand{\arraystretch}{1.25}
\begin{tabularx}{\linewidth}{@{}>{\ttfamily\footnotesize}l c c L{0.40\linewidth}@{}}
\toprule
\sffamily 파일 & \sffamily 전 & \sffamily 후 & \sffamily 변화 \\
\midrule
Web/simulation.js   & $\sim$2,300 & 383   & 거대 클로저 → 얇은 오케스트레이터 \\
Web/Sim\_Parts/ (신규) & ---     & 14모듈 & 기능별 서브시스템 분리 \\
Web/Simulation/ (신규) & ---     & 5모듈  & 토픽별 래퍼$\cdot$오케스트레이터 \\
Web/main.js         & 1,051   & 1,473 & 모바일 내비$\cdot$버튼 동선 보강 \\
Web/ui.js (신규)     & ---     & 97    & \code{main.js}에서 UI 보조 분리 \\
Web/blocklyconfig.js & ---    & {\footnotesize $\pm$135} & 속도 인자$\cdot$탭 이름 블록 \\
Web/bluetooth.js    & ---     & {\footnotesize $\pm$118} & \code{GET\_NAMES} 등 프로토콜 \\
Web/commandexecutor.js & ---  & {\footnotesize $\pm$56}  & 명령에 속도 인자 반영 \\
Web/styles.css      & ---     & {\footnotesize $+$455}   & 모바일 하단 내비$\cdot$반응형 \\
Pico/system.json    & 존재    & 삭제  & JSON 설정 폐지 \\
Pico/ModelFactorySetting.txt (신규) & --- & 55 & 텍스트(키=값) 설정으로 대체 \\
Pico/system\_config.py & ---  & {\footnotesize $\pm$305} & 텍스트 파서$\cdot$주석 보존 저장 \\
Pico/process\_data.py & 653   & 713   & 속도 기능$\cdot$모터 속도 제한 \\
\bottomrule
\end{tabularx}
\end{table}

\noindent 이하 D.2$\sim$D.5는 부록 A의 해당 장을 기준으로 \emph{무엇이 어떻게 달라졌는지}를
트랙별로 정리한다.

\subsection{시뮬레이터 --- 모놀리스 \code{simulation.js}의 모듈화 (유형 A 실행)}
\textbf{[부록 A 기준]} 부록 A(\S A.2 \code{simulation.js}, 8장 미션 시뮬레이터)는
\code{simulation.js}를 three.js 기반 3D 시뮬레이션의 \emph{초기화$\cdot$관리} 단일
파일로 기술한다. 공개 API는 \code{setupSimulation(\{workspace, onOpen, onClose\})}
하나뿐이고, \code{OLED\_ICONS}$\cdot$\code{TOPICS}$\cdot$\code{buildSim}$\cdot$%
\code{recolorLaunchpadAntenna} 등은 \emph{거대 클로저 \code{buildSim()} 내부의 비공개
헬퍼}였다(1차 진단의 ``유형 A --- 거대 클로저'').

\smallskip
\noindent\textbf{[통합 후]} 「2차 Refactoring 계획」의 유형 A 실행으로,
\code{buildSim()}에 묶여 있던 약 50개 상태와 서브시스템을 끌어올려 \emph{공유
컨텍스트(ctx) + 기능별 모듈} 구조로 분해했다. \code{simulation.js}는 \textbf{383줄의 얇은
오케스트레이터}로 줄고, 기능은 \code{Web/Sim\_Parts/} 14개 모듈과 \code{Web/Simulation/}
5개 래퍼로 분리됐다(\code{setupSimulation()} 외부 인터페이스는 그대로라 \code{main.js}는
변경 거의 없음).

\begin{table}[H]\centering\small
\caption{거대 클로저 \code{buildSim()} → 모듈 분해 (유형 A)}
\renewcommand{\arraystretch}{1.2}
\begin{tabularx}{\linewidth}{@{}>{\sffamily\bfseries}l L{0.62\linewidth}@{}}
\toprule
\sffamily 모듈(\code{Web/Sim\_Parts/}) & \sffamily 책임 \\
\midrule
context.js  & 공유 상태(\code{SimContext}) + 씬$\cdot$카메라$\cdot$렌더러 구성, 서브시스템 조립 \\
render.js / assets.js & 렌더 루프 / GLTF 로딩 \\
movement.js / leds.js / gun.js / oled.js & 주행$\cdot$LED$\cdot$발사$\cdot$디스플레이 동작 \\
rocket.js / traffic.js / waves.js & 로켓 발사$\cdot$신호등$\cdot$음파 연출 \\
audio.js / dispatch.js / topics.js / editor\_controls.js & 음향$\cdot$명령 디스패치$\cdot$토픽 정의$\cdot$객체 배치 편집 \\
\bottomrule
\end{tabularx}
\end{table}

\noindent \code{Web/Simulation/}(\code{Simulation\_Main$\cdot$Launcher$\cdot$Rover$\cdot$%
Traffic$\cdot$AresRobot})은 위 라이브러리를 토픽별로 감싸 \code{buildSim}을 재노출하는
얇은 래퍼다. 머지에는 음파 생성(\code{SpawnSoundWave}) 보정 등 소규모 버그 수정과,
모바일 객체 배치(Move$\cdot$Rotate$\cdot$Scale)를 위한 \code{editor\_controls.js}가 포함된다.
구 모놀리스 전문은 \code{Web/simulation\_backup.js}로 보존된다.

\subsection{피코 펌웨어 --- 설정 체계와 명령 프로토콜 변경}
\textbf{[부록 A 기준]} 부록 A(2장 구성요소, 4장 설정 로딩, \S A.1 \code{system\_config.py})는
설정 \emph{정본}을 \code{system.json}(JSON)으로 기술한다 --- \code{CONFIG\_FILE='system.json'},
모듈 활성 플래그$\cdot$핀 할당$\cdot$보정값을 저장하고, \code{SystemConfig}가 입출력을
담당하며, 구형 \code{calibration.json}을 병합한다. 또한 9장은 설정이 펌웨어
\code{system.json}과 웹 \code{localStorage}에 \emph{이중으로} 존재함을 지적했다.

\smallskip
\noindent\textbf{[통합 후]} JSON 설정을 폐지하고, 교사$\cdot$어린이가 메모장으로 직접 고치는
\textbf{텍스트 설정 \code{ModelFactorySetting.txt}}로 대체했다.
\begin{itemize}
  \item \textbf{설정 파일 텍스트화} --- \code{system\_config.py}에 \code{키=값}(또는
  \code{키:값}) 파서를 추가, 첫 부팅 시 설명이 담긴 기본 템플릿을 자동 생성하고,
  저장 시 \emph{기존 주석$\cdot$가이드라인을 보존}하고 값만 갱신. \code{system.json}은 삭제.
  \item \textbf{속도 기능 + 모터 속도 제한} --- 서보$\cdot$DC 이동 명령이 속도 인자를 받도록
  형식을 확장하고(아래 표), \code{max\_speed} 기준으로 스케일링해 \textbf{0\,--\,100\%로
  제한}(5일차 모터 속도 제한 코드). \code{process\_data.py}는 653$\to$713줄.
  \item \textbf{커스터마이징 확장} --- 블록코딩 탭 한글 이름(\code{GET\_NAMES}), 테마 모델
  (로버/발사대), 모듈 토글, GP 핀 번호, 서보 좌우 보정치를 모두 설정 파일에서 제어.
\end{itemize}

\begin{table}[H]\centering\small
\caption{명령 프로토콜 변경 --- 부록~\ref{app:code} 명령 목록 대비 신규$\cdot$확장}
\renewcommand{\arraystretch}{1.2}
\begin{tabularx}{\linewidth}{@{}>{\ttfamily\footnotesize}l L{0.54\linewidth}@{}}
\toprule
\sffamily 명령 형식 & \sffamily 의미(신규/확장) \\
\midrule
SERVO\_방향,속도          & 서보 휠 연속 이동 + 속도(0--100\%) \\
SERVO\_t방향,초,속도       & 시간 제한 서보 이동 + 속도 \\
DC\_방향,속도 / DC\_t방향,초,속도 & DC 모터 연속/시간 제한 + 속도 \\
GET\_NAMES                & 블록코딩 탭 한글 이름 조회(신규) \\
\bottomrule
\end{tabularx}
\end{table}

\begin{notebox}[주의 --- 계약 문서 갱신 필요]
명령 프로토콜이 바뀌었으므로(예: \code{SERVO\_방향,속도}, \code{GET\_NAMES}) 부록 A가
사실상의 계약 문서로 지목한 \code{API\_DOCUMENTATION.md}와 펌웨어$\cdot$웹 양측을
\textbf{짝으로} 갱신해야 한다. 기존 기기의 \code{system.json} 값은 새
\code{ModelFactorySetting.txt}로 자동 이전되지 않는다(기본 템플릿 초기화).
\end{notebox}

\subsection{웹 인터페이스 --- 라우팅 정리와 블록코딩$\cdot$BLE 연동}
\textbf{[부록 A 기준]} 부록 A(\S A.2 \code{main.js})는 \code{main.js}를 라우팅$\cdot$뷰
전환$\cdot$Blockly 초기화$\cdot$이벤트$\cdot$버튼 콜백을 \emph{대부분 내부 함수로} 담당하는
진입점으로 기술한다(1차 진단에서 ``한 함수가 미션/저장/예제/AI/대시보드를 전부 처리''로
지목). 7장은 새 블록 추가 시 \code{blocklyconfig.js}$\cdot$툴박스$\cdot$\code{commandexecutor.js}
세 곳을 함께 고쳐야 함을 강조했다.

\smallskip
\noindent\textbf{[통합 후]}
\begin{itemize}
  \item \textbf{UI 보조 분리} --- \code{updateBlockCodingButtonUI}$\cdot$\code{setupLogToggle}$\cdot$%
  \code{setupContentToggle}를 \textbf{\code{ui.js}(신규 97줄)}로 추출(\code{main.js}는
  1,051$\to$1,473줄로, 모바일 동선 로직 추가분 포함).
  \item \textbf{모바일 하단 내비 + 버튼 동선} --- \emph{개요$\cdot$차시$\cdot$미션$\cdot$AI$\cdot$로그}
  하단 고정 바(\code{safe-area-inset} 대응)와 상단 버튼(개요$\cdot$블록코딩$\cdot$점검) 이동
  버그 수정. \code{styles.css}는 모바일$\cdot$반응형으로 $+$455줄.
  \item \textbf{명령 프로토콜 연동} --- 속도 기능$\cdot$\code{GET\_NAMES}에 맞춰
  \code{blocklyconfig.js}$\cdot$\code{bluetooth.js}$\cdot$\code{commandexecutor.js}를 함께 갱신
  (부록 A가 지적한 ``세 곳 동시 수정'' 지점이 그대로 반영됨).
\end{itemize}

\subsection{진입점$\cdot$자산 --- 페이지 구성과 3D 모델 경량화}
\textbf{[부록 A 기준]} 부록 A는 진입점으로 \code{main.html}(블록 편집기)$\cdot$%
\code{dashboard.html}(점검)을, 모바일 미리보기로 \code{mobile-preview.js}를 기술한다.

\smallskip
\noindent\textbf{[통합 후]} 공개 랜딩(루트 \code{index.html})$\cdot$개발 랜딩(\code{dev.html})과
메인 앱(\code{Web/index.html})으로 페이지 구성이 확장됐고(\code{main.html}/\code{styles.css}에
모바일 내비$\cdot$3D 히어로 추가), 3D 모델 \code{ares\_robot.glb}는 \textbf{7.21\,MB $\to$ 0.58\,MB}로
경량화(약 92\% 감소), 우주인 모델 \code{Astronaut.glb}가 추가됐다. 이 코드 변경은
본문(CLAUDE.md ``Web UI Module Structure'')에도 반영했다.

\begin{infobox}[title=중간 점검 결론]
\small
1$\cdot$2차 Refactoring 계획(부록~\ref{app:refactor1}$\cdot$\ref{app:refactor2})의 핵심
항목 --- \emph{유형 A 거대 클로저 분할}(시뮬레이터)과 \emph{피코 설정$\cdot$명령 체계
정비} --- 이 통합 머지로 실제 코드에 반영됐다. 남은 과제는 부록 A가 짚은 \emph{설정
이중화}(\code{system.json}$\leftrightarrow$\code{localStorage})의 단일화 마무리와,
바뀐 명령 프로토콜에 맞춘 \code{API\_DOCUMENTATION.md} 갱신이다.
\end{infobox}
